\section{Exact-cone trials, quotient mass and weighted lower bounds}
\label{sec:model-energy}

This section fixes the central and cap geometries while their scales vary
independently. The estimates are statements about the resulting surgery
metrics. The comparison with the actual Ricci-flat metrics is a separate
step.

\subsection{Fixed models and central representatives}

Fix an outer radius $R>0$ and a cap parameter $A>0$. Let $g_0^R$ be
the central metric made exactly conical on each punctured node region
$0<r<R$, and let $g_i^A$ be the completed smoothing cap, exactly
conical beyond radius $2A$. These are the fixed models of
\cite[Sections~4--5]{wuebben-hodge}. On the $i$-th cone write
$g_0^R=\dd r^2+r^2h_i$. Its unit-period harmonic link two-form
$\vartheta_i$ is normalized by
\begin{equation}\label{eq:link-normalization}
 \|\vartheta_i\|_{L^2(h_i)}^2=c_i=\frac{2\pi}{3}.
\end{equation}
The link is the ordinary conifold link, diffeomorphic to $S^2\times S^3$.

Let $M_0$ denote the punctured central space and identify
$V=H^2(M_0;\RR)$ with $H^2(Y;\RR)$. The fixed central Hodge input is
the following precise statement. Each class $D\in V$ has a central
$L^2$ harmonic representative $\eta_D$, chosen linearly, and on every
end
\begin{equation}\label{eq:central-end-expansion}
 \eta_D=a_i(D)\vartheta_i+\dd v_{i,D},\qquad
 |v_{i,D}|\le C\|D\|r^{\beta-1},\qquad
 |\dd v_{i,D}|\le C\|D\|r^{\beta-2}
\end{equation}
for a common $\beta>0$, with the corresponding differentiated estimates.
The coefficients $a_i(D)$ are the actual link restrictions
$D\cdot C_i$. This is the fixed-model harmonic expansion proved in
\cite[Lemma~5.1]{wuebben-hodge}. Its finite-dimensional representatives
and primitives are fixed before the cap scales vary.

Define the central mass and surviving subspace by
\begin{equation}\label{eq:model-central-mass}
 G_R(D,F)=\langle\eta_D,\eta_F\rangle_{g_0^R},\qquad
 r(D)=(a_i(D))_i,\qquad K=\ker r,\qquad W_R=K^{\perp_{G_R}}.
\end{equation}
The restriction of $G_R$ to $W_R$ is denoted $M_R$. The attachment
calculation identifies the degree-two cohomology of every smooth
surgery with $K$. For our family, $\dim V=4$ and
$\dim K=\dim W_R=2$.

For arbitrary scales $\varepsilon_i>0$ with
$\ell_i=2A\varepsilon_i<R$, glue the scaled metric
$\varepsilon_i^2g_i^A$ to the central part along its exact cone.
The resulting closed smooth manifold and metric are denoted
$(N_\varepsilon,g_\varepsilon)$. Set
\begin{equation}\label{eq:model-lengths-energy}
 L_i=\log(R/\ell_i),\qquad L_*=\min_iL_i,\qquad
 E_\varepsilon(D,F)
     =\sum_i\frac{c_i}{L_i}a_i(D)a_i(F).
\end{equation}
Every assertion in this section is uniform as $L_*\to\infty$;
no ratio of lengths is bounded. Constants may depend on the fixed
$R,A$ and central and cap models.

\subsection{Trials with exact energy}

Let $\chi_i$ vanish for $r\le\ell_i$, equal
$\log(r/\ell_i)/L_i$ on $\ell_i<r<R$, and equal one for $r\ge R$.
Define a global two-form $T_\varepsilon D$ to be $\eta_D$ outside
the ends, zero on the caps $r\le\ell_i$, and
\begin{equation}\label{eq:corrected-trial}
 T_\varepsilon D
   =a_i(D)\chi_i\vartheta_i+\dd(\chi_i v_{i,D})
\end{equation}
on the conical annulus.

The differential here is distributional. The normal component may jump
at an interface, since $\dd\chi_i$ does. Its tangential trace is zero
at the inner interface and agrees with $\eta_D$ at the outer one.
There is consequently no delta contribution to its exterior derivative.
Equivalently the exact correction is differentiated from a continuous
piecewise $H^1$ one-form on each collar. This gives a trial in the graph
domain of $\dd$, without yet asserting membership in the graph domain
of $\delta_\varepsilon$.

Let $Z_\varepsilon$ be the closed subspace of distributionally closed
$L^2$ two-forms on $N_\varepsilon$, and let $P_\varepsilon$ be
orthogonal projection onto $Z_\varepsilon^\perp$.

\begin{lemma}\label{lem:exact-trial-energy}
The coexact trial $U_\varepsilon D=P_\varepsilon T_\varepsilon D$
belongs to the Hodge form domain and satisfies
\begin{equation}\label{eq:exact-model-energy}
 q_\varepsilon(U_\varepsilon D)
 =\|\dd T_\varepsilon D\|^2
 =\sum_i\frac{c_i|a_i(D)|^2}{L_i},
\end{equation}
where $q_\varepsilon(u)=\|\dd u\|^2+\|\delta_\varepsilon u\|^2$.
Under central identification and extension by zero over missing tips,
\begin{equation}\label{eq:unprojected-trial-convergence}
 \|\widetilde T_\varepsilon D-\eta_D\|_{L^2(M_0,g_0^R)}
       \le C L_*^{-1}\|D\|_{G_R}.
\end{equation}
\end{lemma}
\begin{proof}
The exact correction contributes no exterior derivative. Thus
$\dd T_\varepsilon D=\sum_i a_i(D)\dd\chi_i\wedge\vartheta_i$,
including the interfaces just described. On a six-dimensional cone,
$|\dd r\wedge\vartheta_i|^2=r^{-4}|\vartheta_i|_{h_i}^2$ and
$\dd V=r^5\dd r\,\dd V_{h_i}$. The energy on the $i$-th annulus is
\[
 c_i|a_i(D)|^2\int_{\ell_i}^R r|\chi'_i|^2\,\dd r
   =\frac{c_i|a_i(D)|^2}{L_i^2}
          \int_{\ell_i}^R\frac{\dd r}{r}
   =\frac{c_i|a_i(D)|^2}{L_i}.
\]
The annuli are disjoint, which proves the unprojected identity.

The removed part $(I-P_\varepsilon)T_\varepsilon D$ is closed.
Hence $\dd U_\varepsilon D=\dd T_\varepsilon D$. Orthogonality
to every smooth exact two-form gives
$\delta_\varepsilon U_\varepsilon D=0$ distributionally. On a
closed smooth manifold the joint graph domain of $\dd,\delta$ is
the Hodge form domain, and $Z_\varepsilon^\perp$ is its coexact
subspace. This proves \eqref{eq:exact-model-energy}.

On an annulus the central difference is
\[
 a_i(D)(\chi_i-1)\vartheta_i+
       (\chi_i-1)\dd v_{i,D}+\dd\chi_i\wedge v_{i,D}.
\]
Its squared link-term norm, including the missing tip, is bounded by
a constant times $\|D\|^2$ times
\[
 \ell_i^2+L_i^{-2}\int_0^R r\log^2(R/r)\,\dd r.
\]
The remaining squared norm is bounded by a constant times $\|D\|^2$
times
\[
 \ell_i^{2\beta+2}
   +L_i^{-2}\int_0^R r^{2\beta+1}
                     (1+\log^2(R/r))\,\dd r.
\]
Both integrals are finite. Since $\ell_i=Re^{-L_i}$, the tip terms
are also bounded by a constant times $L_*^{-2}$. Summing over
finitely many ends proves \eqref{eq:unprojected-trial-convergence}.
\end{proof}

An end on which $a_i(D)=0$ contributes exactly zero to the energy.
This is essential for relative estimates. Cutting off the whole
representative without the exact correction would introduce an
unwanted additive term on such an end.

\subsection{The mass after coexact projection}

The strong convergence of $T_\varepsilon D$ does not alone control
the norm of $P_\varepsilon T_\varepsilon D$. To do so we construct
closed compactly supported dual forms.

\begin{lemma}\label{lem:compact-duals}
For every sufficiently small $\rho>0$ there is a linear map
$D\mapsto\Psi_{D,\rho}$ from $V$ to smooth compactly supported
closed four-forms on $M_0$, vanishing on every end for $r<\rho$,
such that
\begin{equation}\label{eq:compact-dual-properties}
 \|\Psi_{D,\rho}-*_0\eta_D\|_{L^2(g_0^R)}
       \le C\rho\|D\|_{G_R},\qquad
 \int_{M_0}\eta_F\wedge\Psi_{D,\rho}=G_R(F,D).
\end{equation}
The constant is independent of $\rho$.
\end{lemma}
\begin{proof}
The form $\omega_D=*_0\eta_D$ is closed and has size at most
$C\|D\|r^{-2}$. Write it as
$\alpha_D(r)+\dd r\wedge\beta_D(r)$ in fixed link coordinates.
Then
$|\alpha_D(r)|_h\le C\|D\|r^2$ and
$|\beta_D(r)|_h\le C\|D\|r$.
Closedness gives
$\dd_L\alpha_D=0$, $\partial_r\alpha_D=\dd_L\beta_D$.
The tangential three-form
$B_D(r)=\int_0^r\beta_D(s)\,\dd s$ converges with its link
derivatives and satisfies
$\dd B_D=\omega_D$, since $\alpha_D(0)=0$. Its conic norm is
at most $C\|D\|r^{-1}$.

Choose $\chi_\rho=1$ for $r\le\rho$, zero for $r\ge2\rho$, with
$|\dd\chi_\rho|\le C/\rho$, and set
\[
 \Psi_{D,\rho}=*_0\eta_D-\sum_i\dd(\chi_\rho B_{i,D}).
\]
This is closed and vanishes near every tip. The squared norm of its
change is bounded by $C\|D\|^2\int_0^{2\rho}r\,\dd r$, including
the cutoff-derivative term, which has size $O(\rho^{-2})$ on its
shell. This proves the first estimate.

For the pairing, integrate
$\eta_F\wedge\dd(\chi_\rho B_D)$ on the space truncated at $r=\delta$.
The tip boundary term is bounded by
$C\|F\|\|D\|\delta^5\delta^{-2}\delta^{-1}
=C\|F\|\|D\|\delta^2$. It tends to zero. All volume integrals
converge, so Stokes makes the correction pairing exactly zero.
\end{proof}

\begin{proposition}\label{prop:trial-quotient-mass}
For $D\in W_R$ and sufficiently large $L_*$,
\begin{equation}\label{eq:projection-mass-loss}
 \|(I-P_\varepsilon)T_\varepsilon D\|
       \le C L_*^{-1}\|D\|_{G_R}.
\end{equation}
For $D,F\in W_R$,
\begin{equation}\label{eq:trial-gram}
 \bigl|\langle U_\varepsilon D,U_\varepsilon F\rangle
                         -G_R(D,F)\bigr|
       \le C L_*^{-1}\|D\|_{G_R}\|F\|_{G_R},\qquad
 q_\varepsilon(U_\varepsilon D,U_\varepsilon F)
       =E_\varepsilon(D,F).
\end{equation}
Consequently, if $\nu_j(\varepsilon)$ are the generalized
eigenvalues of $(E_\varepsilon,M_R)$ on $W_R$, the ordered positive
coexact two-form eigenvalues satisfy
\begin{equation}\label{eq:model-upper-bounds}
 \lambda^{\mathrm{coex}}_{2,j}(N_\varepsilon)
       \le(1+C/L_*)\nu_j(\varepsilon),\qquad
                   1\le j\le\dim W_R.
\end{equation}
\end{proposition}
\begin{proof}
If $\max_i\ell_i<\rho$, extend $\Psi_{D,\rho}$ by zero onto
$N_\varepsilon$. For a smooth closed two-form $z$ on that space,
its exterior class lies in $K$; call it $k$. Compactly supported
cohomological pairing gives
\[
 \int_{N_\varepsilon}z\wedge\Psi_{D,\rho}
     =\int_{M_0}\eta_k\wedge\Psi_{D,\rho}
     =G_R(k,D)=0.
\]
Only the class on a common exterior containing the support is used;
no smooth extension of $z$ over the missing central tips is required.
The equality extends to every closed $L^2$ form by heat regularization
on each fixed closed smoothing. This regularization needs no estimate
uniform in $\varepsilon$.

Since $*^2=1$ in degrees two and four in dimension six,
$V_{\varepsilon,\rho}D=*_\varepsilon\Psi_{D,\rho}$ is perpendicular
to $Z_\varepsilon$. Its support is in the unchanged central metric.
Lemma~\ref{lem:compact-duals} and
\eqref{eq:unprojected-trial-convergence} give
\[
 \|T_\varepsilon D-V_{\varepsilon,\rho}D\|
       \le C(L_*^{-1}+\rho)\|D\|_{G_R}.
\]
Take $\rho=R/(1+L_*)$. For large $L_*$, the support condition and
$2\rho<R$ hold. Orthogonal projection now proves
\eqref{eq:projection-mass-loss}.

The unprojected Gram matrix converges at rate $O(L_*^{-1})$ by
\eqref{eq:unprojected-trial-convergence}. The projected Gram matrix
differs by
\[
 -\langle(I-P_\varepsilon)T_\varepsilon D,
                (I-P_\varepsilon)T_\varepsilon F\rangle,
\]
which is $O(L_*^{-2})\|D\|\|F\|$. This proves the mass assertion;
polarization of \eqref{eq:exact-model-energy} proves the energy
identity. On $W_R$, the energy is positive definite. The mass bound
makes $U_\varepsilon$ injective. Apply coexact min--max to the first
$j$ generalized eigendirections to obtain
$\nu_j/(1-C/L_*)$, which implies \eqref{eq:model-upper-bounds}.
\end{proof}

The proof uses orthogonality to the full surviving space $K$.
Orthogonality only to the polarizing class would not annihilate every
closed smoothing form in the compactly supported pairing.

\subsection{A fixed exterior projection}

Let $C$ be the compact central exterior with boundary $r=r_*=R/2$.
It has fixed geometry and $H^2(C;\RR)=V$. Let $Z_C$ be the maximal
distributional closed subspace of $L^2$ two-forms on $C$, imposing
no tangential trace condition, and let $\Pi_C$ be its orthogonal
projection. Absolute harmonic representation gives a bounded class
map from $Z_C$ to $V$. Set
\begin{equation}\label{eq:model-observer}
 \cA_\varepsilon u
   =\pr_{W_R}^{G_R}[\Pi_C(u|_C)].
\end{equation}
The subscript specifies the changing domain; the exterior operation
itself is fixed.

\begin{lemma}\label{lem:exterior-projection}
For $u|_C$ in the maximal domain of $\dd$,
\begin{equation}\label{eq:exterior-derivative-bound}
 \|(I-\Pi_C)u\|_{H^1(C)}\le C\|\dd u\|_{L^2(C)}.
\end{equation}
\end{lemma}
\begin{proof}
Put $v=(I-\Pi_C)u$. Then $\dd v=\dd u$. For every one-form
$\phi$ in the maximal domain of $\dd$, its differential lies in
$Z_C$, so $\langle v,\dd\phi\rangle=0$. By the adjoint-domain
definition, $v$ belongs to the minimal codifferential domain and
$\delta v=0$. This includes its weak zero normal trace.

The first-order absolute boundary regularity theorem gives
\[
 \|v\|_{H^1(C)}
       \le C(\|\dd v\|+\|\delta v\|+\|v\|).
\]
The weak domain and regularity used here are those of
\cite[Definition~1.2.5, Lemma~1.2.10 and Theorem~4.3.8]{frey}.
There is no imposed normal trace condition on $\dd v$.

On this subspace also $\|v\|\le C\|\dd v\|$. Otherwise unit
vectors in $Z_C^\perp$ with differential tending to zero would be
$H^1$ bounded. Rellich compactness would give a unit strong $L^2$
limit in both $Z_C$ and $Z_C^\perp$, a contradiction. Combining
the two estimates proves \eqref{eq:exterior-derivative-bound}.
\end{proof}

\subsection{Boundary coefficients and capacity}

For a smooth two-form, let
\begin{equation}\label{eq:boundary-coefficient}
 B_i(u)=c_i^{-1}\int_{\{r=r_*\}}u_T\wedge *_h\vartheta_i.
\end{equation}
This integrates the tangential trace over the whole five-dimensional
boundary against a smooth three-form. It is a bounded $H^1$ trace
functional, not an integral of an arbitrary $H^1$ two-form over a
codimension-three cycle.

For $\zeta=\Pi_Cu$, the harmonic boundary coefficient is the
restriction coordinate $a_i([\zeta])$. This follows by pairing the
closed tangential class with $*_h\vartheta_i$. For smooth data,
first-order boundary regularity makes $\zeta$ smooth; approximation
then gives the assertion in the weak setting. The fixed trace
estimate and Lemma~\ref{lem:exterior-projection} yield
\begin{equation}\label{eq:boundary-observer-error}
 |B_i(u)-a_i(\cA_\varepsilon u)|
       \le C\|\dd u\|_{L^2(C)}
       \le C\sqrt{q_\varepsilon(u)}.
\end{equation}
Here $r$ vanishes on $K$, so taking the quotient projection has not
changed the restriction coefficient.

The cap and annulus capacity bound, with its exact normalization, is
\begin{equation}\label{eq:cap-capacity}
 q_{\varepsilon,i}(u)
       \ge\frac{c_i|B_i(u)|^2}
                    {\log(r_*/\varepsilon_i)+C_i}.
\end{equation}
We recall its variational reason from
\cite[Section~4]{wuebben-hodge}. The closed three-form
$*_h\vartheta_i$ extends into the fixed smoothing cap. Stokes and
Cauchy--Schwarz bound its inner boundary coefficient by the cap
$\dd$-energy. On the conical annulus the harmonic component of
$\dd u$ controls $c_i\int r|B'_i(r)|^2\,\dd r$. Adding the two
resistances gives \eqref{eq:cap-capacity}. The energy of a two-form
in real dimension six is invariant under constant metric scaling,
so the cap constant is independent of $\varepsilon_i$. No special
boundary condition is imposed on the competitor, and no other cap
scale enters.

\begin{proposition}\label{prop:weighted-observer-bound}
For every form in the Hodge form domain, and sufficiently large $L_*$,
\begin{equation}\label{eq:model-weighted-lower}
 E_\varepsilon(\cA_\varepsilon u)
       \le(1+CL_*^{-1/2})q_\varepsilon(u).
\end{equation}
Moreover,
\begin{equation}\label{eq:observer-bounded-closed}
 \begin{aligned}
 \|\cA_\varepsilon u\|_{G_R}\le C\|u\|,\qquad
 &\cA_\varepsilon z=0\quad(z\in Z_\varepsilon),\\
 \|\cA_\varepsilon U_\varepsilon D-D\|_{G_R}
       &\le CL_*^{-1}\|D\|_{G_R}\quad(D\in W_R).
 \end{aligned}
\end{equation}
\end{proposition}
\begin{proof}
Put $w_i=c_i/L_i$. The denominator in \eqref{eq:cap-capacity}
is $L_i$ plus a fixed bounded shift. Sum over the disjoint regions
to obtain
\[
 \sum_iw_i|B_i(u)|^2\le(1+C/L_*)q_\varepsilon(u).
\]
Equation~\eqref{eq:boundary-observer-error}, the fixed number of
ends and $w_i\le C/L_*$ give
\[
 \sum_iw_i|B_i(u)-a_i(\cA_\varepsilon u)|^2
       \le C L_*^{-1}q_\varepsilon(u).
\]
The triangle inequality in this weighted Euclidean norm proves
\eqref{eq:model-weighted-lower}, after squaring. Every boundary
error is proportional to the energy of the same competitor.
Smooth density extends the estimate to the form domain.

For boundedness of $\cA_\varepsilon$, restriction and $\Pi_C$ are
contractions, and absolute harmonic representation and the fixed
finite quotient projection are bounded. A globally closed form has
exterior class in $K$, proving the second assertion; closed $L^2$
forms follow by regularization on each smoothing. Finally
$\cA_\varepsilon\eta_D=D$ on $W_R$ and
\eqref{eq:unprojected-trial-convergence} gives the last estimate
for $T_\varepsilon D$. Its difference from $U_\varepsilon D$ is
globally closed, so both have the same observation.
\end{proof}

The trials and weighted observation now verify the first two hypotheses
of Theorem~\ref{thm:relative-reduction} on the model coexact sector.
The remaining mass-capture estimate concerns arbitrary low-energy forms.
It requires independent-cap compactness, not merely the mass convergence
of the trials proved above.
